\section{Алгоритм с изменяемой длительностью измерительного окна, кратной периоду инжектируемого сигнала}\label{sec:fr_optim}

\needspace{1\baselineskip}
{Принципиальные отличия рассматриваемого алгоритма от алгоритма \ref{sec:Th_const}} заключаются в том, что длительность измерительного окна ${T_h}$, длина последовательностей $N$ и период дискретизации ${T_s}$ -- плавающие параметры, соответствующие каждому отдельному значению периода инжектируемого сигнала ${T_z}$. Остаются следующие ограничения:
\begin{itemize}
    \item Возможные значения длительности измерительного окна ${T_h}$ кратны периоду инжектируемого сигнала ${T_z}$.

    \item Период дискретизации ${T_s}$ кратен периоду выполнения операций в системе ${T_d}$.
\end{itemize}
Эти ограничения сводятся к плавающим параметрам ${L, M \in \mathbb{N}_{>0}}$, в дополнение к плавающему параметру ${N \in \mathbb{N}_{>0}}$.

Таким образом, в серии из $K$ измерений для каждого значения периода инжектируемого сигнала ${T^{(k)}_z}$, ${k \in \{1, \dots, K\}}$ (рисунок~\ref{fig:Th_var_N_var}), выполняется:
\begin{equation}\label{eqn:windowlength}
    T^{(k)}_h = L^{(k)} \mathord{\cdot} T^{(k)}_z = N^{(k)} \mathord{\cdot} \underbrace{M^{(k)} \mathord{\cdot} T_d}_{T^{(k)}_s} = L^{(k)} \mathord{\cdot} R^{(k)} \mathord{\cdot} T_d,
\end{equation}
где ${R^{(k)} \in \mathbb{R}_{>0}}$ -- отношение периода инжектируемого сигнала ${T^{(k)}_z}$ к периоду выполнения операций в системе ${T_d}$.

\begin{figure}[H]
    \centering
    \includegraphics [width=1.00\linewidth] {frequencyResponse/Th_var_N_var.pdf}

    \caption{Параметры в алгоритме с изменяемой длительностью измерительного окна, кратной периоду инжектируемого сигнала}
    \label{fig:Th_var_N_var}
\end{figure}

\textbf{Аналогично алгоритму \ref{sec:Th_const}}, соотношение между периодом инжектируемого сигнала и периодом дискретизации может принимать дробные значения. В совокупности с изменяемой длительностью измерительного окна это обеспечивает на интервале ${f \in \left(0; {f_d / 2} \right)}$ сколь угодно много потенциальных значений частот ${f_z}$ инжектируемого сигнала, что \textbf{является главной отличительной особенностью алгоритма}.

Из равенства \eqref{eqn:windowlength}
\begin{equation}\label{eqn:L_ratio}
    L^{(k)} = \frac{N^{(k)} \mathord{\cdot} M^{(k)}}{R^{(k)}}.
\end{equation}

В связи с тем, что частота инжектируемого сигнала кратна частоте первой гармоники в спектре:
\begin{equation*}
    f^{(k)}_z = L^{(k)} \mathord{\cdot} f^{(k)}_h = \cfrac{L^{(k)}}{T^{(k)}_h},
\end{equation*}
и ограничена половиной частоты дискретизации:
\begin{equation*}
    f^{(k)}_z < \cfrac{f^{(k)}_s}{2} \quad \rightarrow \quad \cfrac{L^{(k)}}{T^{(k)}_h} < \cfrac{f_d}{2 M^{(k)}},
\end{equation*}
множитель ${L^{(k)}}$ ограничен диапазоном
\begin{equation*}
    1 \leq L^{(k)} < N^{(k)} / 2,
\end{equation*}
что, с учетом \eqref{eqn:L_ratio}, накладывает ограничение на коэффициент децимации
\begin{equation*}
    1 \leq M^{(k)} < R^{(k)} / 2.
\end{equation*}

Таким образом, для каждой $k$-й частоты ${f_z}$ задача сводится к поиску наборов ${\mathbf{s}_{(p)} \in \mathbb{N}_{> 0}^3}$ натуральных чисел ${N_{(p)}, M_{(p)}, L_{(p)}}$, при которых частота инжектируемого сигнала ${f_{z(p)}}$
\begin{equation*}
    f_{z(p)} = \cfrac{L_{(p)}}{N_{(p)} \mathord{\cdot} M_{(p)}} f_d
\end{equation*}
отклоняется от желаемой частоты ${f_z}$ в пределах заданного значения ${\Delta\epsilon^{*}_{f}}$:
\begin{equation*}
    \left| f_{z(p)} - f_z \right| = \Delta\epsilon_{f(p)} \leq \Delta\epsilon^{*}_f,
\end{equation*}
при ограничениях
\begin{equation*}
        N_\text{min} \leq N_{(p)} \leq N_\text{max}, \quad 1 \leq M_{(p)} < R/2.
\end{equation*}

Оптимальным считается такой набор ${\mathbf{s}_{(p^{*})}}$, в котором параметр ${N^{*}}$ минимален среди всех наборов ${\mathbf{s}_{(p)}}$, отвечающих условию
\begin{equation}\label{eqn:freqtol}
    \Delta\epsilon_{f(p)} \leq \Delta\epsilon^{*}_{f}.
\end{equation}
В случае, когда ни один набор параметров ${\mathbf{s}_{(p)}}$ не удовлетворяет условию \eqref{eqn:freqtol}, оптимальным считается набор, соответствующий минимальному отклонению ${\Delta\epsilon_{f(p)}}$. Таким образом, согласно выражению \eqref{eqn:meas_duration}, минимум длины массивов-откликов $N$ обеспечивает сокращение продолжительности измерения частотной характеристики.